%% Drafted from the mapped corpus by scripts/draft_topics.py.
% Model: gpt-5.6-luna; source characters: 23156.
\section{Функционално програмиране}

\subsection{Обща характеристика на функционалния стил}

Функционалното програмиране е стил на програмиране, в който основните логически единици са функциите. Функцията е програмна единица, която прилага определено преобразуване върху аргументи и връща резултат. Програмата се изгражда чрез дефиниране, използване и комбиниране на функции.

Основното действие във функционалната програма е обръщението към функция. Основният начин за разделяне на програмата на части е дефинирането на име и свързването му с израз, който изчислява стойността на функцията. Основният начин за комбиниране е суперпозицията на функции: резултатът от една функция се подава като аргумент на друга.

Функциите във функционалния стил могат:

\begin{itemize}
    \item да имат параметри;
    \item да се прилагат върху аргументи;
    \item да приемат други функции като аргументи;
    \item да връщат функции като резултат;
    \item да бъдат дефинирани рекурсивно, тоест чрез обръщения към самите себе си.
\end{itemize}

Функционалният стил е представител на декларативния стил на програмиране. При декларативния подход програмата описва какво трябва да бъде полученият резултат, вместо подробно да задава последователност от промени в паметта. Във функционалната програма обикновено липсват традиционни процедурни елементи като присвояване, цикли, заделяне и промяна на клетки в паметта, както и управление на изпълнението чрез прескачане.

Програмите във функционалния стил могат да се разглеждат като списък от дефиниции на функции или като списък от равенства, които описват зависимостите между стойностите. Изпълнението се свежда до оценяване на изрази и прилагане на функции.

\begin{defn}
\emph{Функционално програмиране} е стил на програмиране, при който програмите се изграждат основно чрез дефиниране, прилагане и композиция на функции, а изчисляването се представя като оценяване на изрази.
\end{defn}

Съществен компонент на функционалната програма е липсата на странични ефекти в чистия функционален език. Когато изразът е чист и се оценява в един и същ контекст, той дава една и съща стойност. Това позволява програмата да се разглежда като описание на зависимости между стойности, а не като поредица от изменения на общо състояние.

\subsection{Изрази и стойности в Scheme}

За илюстрация на функционалните понятия ще бъде използван синтаксисът на Scheme. Изразът е атом или комбинация от изрази.

Атомарните изрази не се декомпозират повече. Типични атомарни изрази в Scheme са:

\begin{itemize}
    \item булевите константи \texttt{\#t} и \texttt{\#f};
    \item числовите константи, например \texttt{15}, \texttt{2/3} и \texttt{-1.23};
    \item знаковите константи, например \texttt{\#\textbackslash a} и \texttt{\#\textbackslash newline};
    \item низовите константи, например \texttt{"Scheme"} и \texttt{"hi"};
    \item идентификаторите, тоест имената, които могат да бъдат свързани със стойности;
    \item вградените функции, например \texttt{+}, \texttt{-}, \texttt{*}, \texttt{/} и \texttt{sqrt}.
\end{itemize}

Комбинация в Scheme е израз, конструиран като списък от изрази, заградени в скоби:

\begin{verbatim}
(<expression0> <expression1> ... <expressionn>)
\end{verbatim}

Първият израз, означен като \texttt{<expression0>}, се нарича оператор. Очаква се неговата оценка да бъде функция. Останалите изрази са операнди. Стойността на комбинацията се получава, като функцията, зададена чрез оператора, се приложи върху стойностите на операндите.

Например:

\begin{verbatim}
(+ 2 3)
\end{verbatim}

е комбинация, чийто оператор е \texttt{+}, а операндите са \texttt{2} и \texttt{3}. Оценката на комбинацията е резултатът от прилагането на функцията за събиране върху двете числа.

Правилата за оценяване на някои основни изрази са следните:

\begin{itemize}
    \item оценката на число е самото число;
    \item оценката на низ е самият низ;
    \item оценката на вграден оператор е реализацията на този оператор;
    \item оценката на име е стойността, свързана с името в текущата среда.
\end{itemize}

\begin{defn}
\emph{Среда} е контекстът, който съдържа свързвания между имена и стойности. При оценяване на име се търси стойността, свързана с това име в текущата среда.
\end{defn}

\subsection{Дефиниране и използване на функции}

Дефинирането на функция представлява свързване на име с израз, който описва изчисляването на резултата. Дефиницията е средство за абстракция, защото позволява сложното изчисление да бъде скрито зад име и използвано многократно.

Дефиниране на име чрез израз в Scheme:

\begin{verbatim}
(define <name> <expression>)
\end{verbatim}

Например:

\begin{verbatim}
(define size 2)
\end{verbatim}

свързва името \texttt{size} със стойността \texttt{2}.

Функция с формални параметри може да се дефинира чрез:

\begin{verbatim}
(define (<name> <formal-parameters>) <body>)
\end{verbatim}

Пример:

\begin{verbatim}
(define (square x)
  (* x x))
\end{verbatim}

Тук \texttt{square} е името на функцията, \texttt{x} е формален параметър, а \texttt{(* x x)} е тялото. При обръщението

\begin{verbatim}
(square 5)
\end{verbatim}

функцията се прилага върху аргумента \texttt{5} и резултатът е \texttt{25}.

\begin{defn}
\emph{Формален параметър} е име, използвано в дефиницията на функцията, което представя стойността на съответния аргумент при конкретно приложение.
\end{defn}

\begin{defn}
\emph{Аргумент} е изразът или стойността, подадени при конкретно обръщение към функция.
\end{defn}

Формалните параметри и действителните аргументи се съпоставят по позиция. При прилагането на функцията формалните параметри получават стойностите на съответните аргументи, след което се оценява тялото на функцията.

Рекурсията е основен метод за дефиниране на функции. При рекурсивна дефиниция функцията се обръща към самата себе си. Източниците посочват рекурсията като характерен начин за дефиниране във функционалното програмиране, без да налагат конкретна рекурсивна задача или пример.

\subsection{Условни форми и локални дефиниции}

Условната форма \texttt{if} има синтаксис:

\begin{verbatim}
(if <condition> <expression1> <expression2>)
\end{verbatim}

Първо се оценява условието. Ако неговата стойност е \texttt{\#t}, се оценява и връща \texttt{<expression1>}. В противен случай се оценява и връща \texttt{<expression2>}.

Тъй като се оценява само избраното разклонение, \texttt{if} е специална форма, а не обикновена функция. Специалните форми са конструкции, които са изключения от общото правило за оценяване на комбинациите.

За многозначен избор се използва \texttt{cond}:

\begin{verbatim}
(cond
  (<condition1> <expression1>)
  ...
  (<conditionn> <expressionn>)
  (else <expressionm>))
\end{verbatim}

Условията се оценяват последователно. При първото условие със стойност \texttt{\#t} се оценява съответният израз и неговата стойност става стойност на цялата форма. Ако нито едно условие не е изпълнено, се избира клонът \texttt{else}.

Локални имена могат да се създават чрез специалната форма \texttt{let}:

\begin{verbatim}
(let ((<symbol1> <expression1>)
      ...
      (<symboln> <expressionn>))
  <body>)
\end{verbatim}

Оценката на \texttt{let} в среда $E$ се описва по следния начин:

\begin{enumerate}
    \item създава се нова среда $E_1$, която е разширение на $E$;
    \item всеки израз $\texttt{<expression}_i\texttt{>}$ се оценява в $E$;
    \item получената стойност се свързва със съответното име $\texttt{<symbol}_i\texttt{>}$ в $E_1$;
    \item тялото се оценява в $E_1$ и неговата стойност се връща като стойност на \texttt{let}.
\end{enumerate}

Символично:

\[
E_1=E\cup
\{\texttt{<symbol}_i\texttt{>} \mapsto
\operatorname{eval}(\texttt{<expression}_i\texttt{>},E)
\mid 1\leq i\leq n\},
\]

а резултатът е:

\[
\operatorname{eval}(\texttt{<body>},E_1).
\]

В Haskell локални дефиниции могат да се задават чрез \texttt{where}:

\begin{verbatim}
<function-definition>
  where
    <definition1>
    ...
    <definitionn>
\end{verbatim}

Областта на действие на тези дефиниции е ограничена до дефиницията на функцията. Дефинициите в \texttt{where} могат да бъдат взаимно рекурсивни.

\section{Оценяване на изрази}

\subsection{Общо правило за оценяване на комбинация}

Общото правило за оценяване на обикновена комбинация се състои от следните стъпки:

\begin{enumerate}
    \item оценяват се подизразите на комбинацията;
    \item оценката на най-левия подизраз се разглежда като функция;
    \item тази функция се прилага върху оценките на останалите подизрази.
\end{enumerate}

Например:

\begin{verbatim}
(* (+ 2 (* 4 6)) (+ 3 5 7))
\end{verbatim}

се оценява чрез последователността:

\begin{verbatim}
(* (+ 2 24) 15)
(* 26 15)
390
\end{verbatim}

Тук общото правило се прилага към вложените комбинации \texttt{(* 4 6)}, \texttt{(+ 2 (* 4 6))}, \texttt{(+ 3 5 7)} и външната комбинация.

\begin{keydefn}
При обикновена комбинация първо се оценяват операторът и операндите, а след това функцията, получена като стойност на оператора, се прилага върху стойностите на операндите.
\end{keydefn}

Конструкциите \texttt{define}, \texttt{if} и \texttt{cond} не следват непосредствено това правило. Те са специални форми и имат собствени правила за оценяване.

\subsection{Оценяване на приложение на функция}

При комбинация, чийто оператор е функция, се оценяват елементите на комбинацията и се прилага процедурата, която е стойност на оператора, върху стойностите на операндите.

При дефиницията:

\begin{verbatim}
(define (square x)
  (* x x))
\end{verbatim}

специалната форма \texttt{define} създава свързване между името \texttt{square} и зададената функция. Това свързване се добавя в текущата среда. Функцията може по-късно да бъде приложена чрез:

\begin{verbatim}
(square 4)
\end{verbatim}

Моделът на заместването представя прилагането на функцията чрез заместване на формалните параметри със съответните аргументи. За \texttt{(square 4)} тялото

\begin{verbatim}
(* x x)
\end{verbatim}

се разглежда след заместването като

\begin{verbatim}
(* 4 4)
\end{verbatim}

и се оценява до \texttt{16}. Оценката на последния израз от тялото е оценката на обръщението към функцията.

Този модел е концептуален: той обяснява как аргументите се свързват с формалните параметри. Реалната реализация може да използва друга вътрешна техника, но поведението се описва чрез същите зависимости.

\section{Модели на оценяване}

При модела на заместването са възможни два основни подхода за оценяване на аргументите: апликативен и нормален.

\subsection{Апликативно оценяване}

Апликативното оценяване се нарича още строго оценяване или оценяване по стойност (\textit{call-by-value}). При него първо се оценяват операндите, след което функцията се прилага върху получените стойности.

Схематично:

\[
\text{оценяване на операндите}
\longrightarrow
\text{получаване на стойности}
\longrightarrow
\text{прилагане на функцията}.
\]

Ако имаме приложение от вида

\begin{verbatim}
(f e1 e2)
\end{verbatim}

апликативният подход първо оценява \texttt{e1} и \texttt{e2}, а след това прилага \texttt{f} върху получените резултати.

Предимството на този подход, посочено в източниците, е, че преди прилагането на функцията всеки аргумент е сведен до една стойност. Това е пестеливо откъм памет, тъй като не се налага да се съхраняват неоценени изрази. Аргументът е, че се извършва оценяване дори когато стойността му в крайна сметка не е необходима на тялото на функцията.

\subsection{Нормално оценяване}

Нормалното оценяване се нарича още лениво оценяване или оценяване по име (\textit{call-by-name}). При него аргументите не се оценяват предварително. Вместо това имената на аргументите се заместват в тялото на функцията и изразите се оценяват, когато това стане необходимо.

Схематично:

\[
\text{заместване в тялото}
\longrightarrow
\text{достигане до примитивни операции}
\longrightarrow
\text{оценяване на необходимите изрази}.
\]

Следователно нормалният подход първо разгъва функцията и чак след това оценява аргументите, които участват в получения израз.

Нормалното оценяване може да избегне оценяване на аргумент, който не се използва. То обаче може да доведе до многократно оценяване на един и същ израз, ако той се появява повече от веднъж след заместването. Така се пести време за предварително оценяване само когато изразът не е необходим, но може да се изразходва повече време при многократно пресмятане. Същевременно е необходимо да се съхраняват или повторно да се разгръщат неоценени изрази, поради което подходът е по-непестелив откъм памет.

\begin{keydefn}
При апликативното оценяване аргументите се оценяват преди прилагането на функцията. При нормалното оценяване аргументите се заместват и се оценяват само когато са необходими.
\end{keydefn}

\begin{remark}
Една и съща програма може да се държи различно при апликативен и нормален подход, особено когато функцията не използва някой аргумент или когато оценяването на аргумент би довело до проблем. Източниците посочват също, че в стандартните езици има лениво поведение в отделни конструкции, например при конюнкция, когато след лъжлив операнд останалите операнди не се пресмятат.
\end{remark}

\begin{supp}[Уточнение за терминологията]
В учебната литература изразът ``лениво оценяване'' често се използва по-широко от ``call-by-name'' и може да включва запомняне на вече изчислена стойност. Разглежданият тук материал представя нормалния подход чрез заместване и акцентира върху възможното многократно оценяване на един и същ израз.
\end{supp}

\section{Функции от по-висок ред}

\subsection{Определение}

\begin{defn}
\emph{Функция от по-висок ред} е функция, която приема функция като параметър или връща функция като резултат.
\end{defn}

Функциите от по-висок ред са естествено следствие от идеята, че функциите са стойности, с които програмата може да работи. Те могат да се подават на други функции, да се връщат като резултат и да се използват за абстрахиране на повтарящи се схеми на изчисление.

Функция, която приема друга функция като параметър, може да прилага подадената функция върху дадена стойност. Например предикатът за неподвижна точка може да бъде записан в Scheme като:

\begin{verbatim}
(define (fixed-point f x)
  (= (f x) x))
\end{verbatim}

Тук \texttt{f} е параметър, чиято стойност трябва да бъде функция, а \texttt{x} е стойност, върху която \texttt{f} се прилага.

Практически примери за функции от по-висок ред са \texttt{map}, \texttt{filter}, \texttt{foldl}, \texttt{foldr} и \texttt{apply}. Функцията \texttt{map} прилага дадена функция върху всеки елемент на списък и връща нов списък:

\begin{verbatim}
(map (lambda (x) (* x x)) '(1 2 3))
\end{verbatim}

Резултатът е:

\begin{verbatim}
'(1 4 9)
\end{verbatim}

Функцията \texttt{filter} приема предикат и списък и връща елементите, за които предикатът е истина:

\begin{verbatim}
(filter even? '(1 2 3 4 5))
\end{verbatim}

Резултатът е:

\begin{verbatim}
'(2 4)
\end{verbatim}

Функциите \texttt{foldr} и \texttt{foldl} свиват списък чрез подадена функция и начална стойност. За пример:

\begin{verbatim}
(foldr - 0 '(1 2 3))
\end{verbatim}

се интерпретира като:

\[
1-(2-(3-0)).
\]

При \texttt{foldl} обработката се извършва отляво надясно чрез акумулиране на междинния резултат:

\begin{verbatim}
(foldl - 0 '(1 2 3))
\end{verbatim}

В Racket функцията за \texttt{foldl} получава първо текущия елемент и после акумулатора, затова изразът се интерпретира като:

\[
3-(2-(1-0))=2.
\]

Този ред на аргументите е различен от често използваната Haskell конвенция, при която акумулаторът е първи.

Функцията \texttt{apply} приема функция и списък от аргументи, като използва елементите на списъка като отделни аргументи:

\begin{verbatim}
(apply + '(1 2 3))
\end{verbatim}

е еквивалентно на:

\begin{verbatim}
(+ 1 2 3)
\end{verbatim}

и дава резултат \texttt{6}.

\subsection{Къринг}

В Haskell функция с няколко аргумента може да се разглежда като функция на един аргумент, която връща функция, очакваща следващия аргумент. Това представяне се нарича къринг.

Ако функцията има тип:

\[
t_1\to t_2\to t_3,
\]

тя приема аргумент от тип $t_1$ и връща функция от тип:

\[
t_2\to t_3.
\]

Следователно частичното прилагане на функцията създава нова функция. Например:

\begin{verbatim}
div50 x = div 50 x
\end{verbatim}

може да се разглежда като дефиниране на:

\begin{verbatim}
div50 = div 50
\end{verbatim}

и след това:

\begin{verbatim}
div50 4
\end{verbatim}

дава \texttt{12}.

\begin{defn}
\emph{Къринг} е представяне на функция с няколко аргумента като последователност от функции, всяка от които приема един аргумент и връща функция или краен резултат.
\end{defn}

\subsection{Анонимни функции}

Анонимна, или ламбда, функция е функция без зададено име. Тя може да бъде конструирана на мястото, където е необходима, особено като аргумент на функция от по-висок ред.

Синтаксисът в Scheme е:

\begin{verbatim}
(lambda (<parameters>) <body>)
\end{verbatim}

Пример:

\begin{verbatim}
(lambda (x) (+ x 3))
\end{verbatim}

оценява до функционален обект с параметър \texttt{x} и тяло \texttt{(+ x 3)}. Този обект може да бъде приложен непосредствено:

\begin{verbatim}
((lambda (x) (+ x 3)) 5)
\end{verbatim}

Резултатът е \texttt{8}.

Анонимните функции позволяват на програмиста да задава поведението на функция от по-висок ред без предварително именуване на помощна функция. В предишния пример ламбда-функцията е подадена директно на \texttt{map}.

\subsection{Функции, които връщат функции}

Функция от по-висок ред може да връща нова функция като резултат. Пример е функцията \texttt{twice}, която приема функция \texttt{f} и връща функция, прилагаща \texttt{f} два пъти:

\begin{verbatim}
(define (twice f)
  (lambda (x) (f (f x))))
\end{verbatim}

Обръщението:

\begin{verbatim}
(twice square)
\end{verbatim}

връща функция. Прилагането на получения резултат върху \texttt{3} дава:

\begin{verbatim}
((twice square) 3)
\end{verbatim}

Това е еквивалентно на прилагане на \texttt{square} два пъти:

\[
(3^2)^2=81.
\]

Друг пример е функция, която създава функция за прибавяне на фиксирано число:

\begin{verbatim}
(define (n+ n)
  (lambda (i) (+ i n)))
\end{verbatim}

Така може да се дефинира:

\begin{verbatim}
(define 1+ (n+ 1))
\end{verbatim}

Полученото \texttt{1+} е функция, която прибавя фиксирана стойност към своя аргумент.

\begin{supp}[Бележка към примерната стойност]
При дефиницията \texttt{(define 1+ (n+ 1))} обръщението \texttt{(1+ 4)} според самото тяло \texttt{(+ i n)} дава \texttt{5}. Среща се и изписване на друга резултатна стойност в учебния материал; тя не следва от показаната дефиниция, затова тук е запазена зависимостта, зададена от кода.
\end{supp}

\section{Изпитен фокус}

\begin{itemize}
    \item Да се дефинира функционалното програмиране и да се опишат функциите като основни програмни единици.
    \item Да се обясни декларативният характер на функционалния стил, композицията на функции и ролята на рекурсията.
    \item Да се различат атом, комбинация, оператор и операнд в Scheme.
    \item Да се възпроизведат синтаксисът и смисълът на \texttt{define}, \texttt{if}, \texttt{cond}, \texttt{let} и локалните дефиниции чрез \texttt{where}.
    \item Да се изложи общото правило за оценяване на комбинация и да се проследи примерът с вложено събиране и умножение.
    \item Да се обясни моделът на заместването при прилагане на функция.
    \item Да се сравнят апликативното оценяване \textit{call-by-value} и нормалното оценяване \textit{call-by-name} по ред на оценяване, време и памет.
    \item Да се дефинира функция от по-висок ред и да се дадат примери с функции като параметри и резултати.
    \item Да се обяснят кърингът, частичното прилагане и анонимните ламбда-функции.
    \item Да се проследи примерът \texttt{twice}, който връща функция, прилагаща подадена функция два пъти.
\end{itemize}
